\section{Arquitectura general}

El Mini Banco Concurrente Distribuido es un sistema de tres nodos que mantiene
en memoria un conjunto de cuentas y atiende cuatro operaciones bancarias por
HTTP. El estado no vive en ningun gestor de base de datos: las cuentas se cargan
desde un CSV al arrancar y residen en estructuras concurrentes del proceso Java.
Esta seccion describe la topologia de los nodos, el modelo de datos, el cableado
del servidor en \codigo{BancoHttp.crear()} y la organizacion de los paquetes.

\subsection{Vision general: tres nodos con lider fijo}

El despliegue usa tres procesos identicos: \codigo{nodo-1} como lider fijo y
\codigo{nodo-2} y \codigo{nodo-3} como replicas. Todas las lecturas y escrituras
de la logica bancaria se dirigen al lider; las replicas siguen al lider y aplican
sus transacciones en vivo, de modo que conservan una copia consistente del estado
sin atender la operacion bancaria directamente.

El rol no se compila ni se configura en codigo: lo decide la variable de entorno
\codigo{BANCO\_LIDER\_HOST}. Si no esta definida, el nodo es LIDER; si esta
definida, el nodo es REPLICA y apunta a ese host. El criterio vive en
\codigo{server/MainServer.java}:

\begin{lstlisting}[language=Java,caption={Eleccion de rol por variable de entorno (MainServer.java)}]
/** LIDER si no hay BANCO_LIDER_HOST; en otro caso REPLICA. */
private static boolean esLider() {
    String liderHost = System.getenv("BANCO_LIDER_HOST");
    return liderHost == null || liderHost.isBlank();
}
\end{lstlisting}

El mismo metodo gobierna el arranque de la replicacion: el lider levanta un
\codigo{ServidorReplicacion} en un puerto TCP, mientras que la replica crea un
\codigo{ClienteReplicacion} que se conecta al lider para recibir y aplicar las
transacciones.

\begin{lstlisting}[language=Java,caption={Arranque de replicacion segun el rol (MainServer.java)}]
if (esLider()) {
    new ServidorReplicacion(CuentaRepository.get(), puertoRepl).iniciar();
    System.out.println("Rol: LIDER - replicacion TCP en puerto " + puertoRepl);
} else {
    String liderHost = System.getenv("BANCO_LIDER_HOST");
    new ClienteReplicacion(CuentaRepository.get(), liderHost, puertoRepl).iniciar();
    System.out.println("Rol: REPLICA - siguiendo a " + liderHost + ":" + puertoRepl);
}
\end{lstlisting}

\begin{nota}
Solo el lider activa el log durable en Cloud Storage (variable
\codigo{BANCO\_GCS\_BUCKET}), y lo hace ANTES de servir trafico para poder
recuperar el estado reaplicando el log sobre las cuentas del CSV en caso de que
todos los nodos hubieran caido.
\end{nota}

\subsection{Modelo de datos en memoria}

El dataset son aproximadamente 820,000 cuentas que se cargan desde
\codigo{alumnos.csv} al iniciar el nodo. La carga la hace
\codigo{data/AlumnosLoader.java}, que lee el archivo linea por linea: el
\codigo{id} de cada cuenta es el numero de linea en orden de lectura empezando en
1. El CSV no tiene cabecera, por lo que la primera fila ya es \codigo{id=1}.

Una decision central del modelo es guardar el dinero SIEMPRE en centavos como
\codigo{long}, nunca como \codigo{double}, para que el invariante de suma
constante de saldos se mantenga exacto. El loader convierte el saldo decimal del
CSV a centavos enteros:

\begin{lstlisting}[language=Java,caption={Carga de cuentas con id base 1 y saldo en centavos (AlumnosLoader.java)}]
id++; // base 1, en orden de lectura
long centavos = new BigDecimal(c[3].trim())
        .movePointRight(2)
        .setScale(0, RoundingMode.HALF_UP)
        .longValueExact();
repo.put(new Cuenta(id, c[0].trim(), c[1].trim(), c[2].trim(), centavos));
\end{lstlisting}

La entidad es \codigo{model/Cuenta.java}. Expone \codigo{id}, los datos del
propietario y el saldo en centavos, y trae un \codigo{ReentrantLock} por cuenta
que las transferencias toman en orden de \codigo{id} para evitar interbloqueos. El
saldo se ofrece como decimal exacto solo al construir el JSON de salida, via
\codigo{saldoDecimal()}, que divide los centavos entre 100 con
\codigo{BigDecimal}.

\begin{lstlisting}[language=Java,caption={Saldo en centavos y conversion a decimal exacto (Cuenta.java)}]
private volatile long saldoCentavos;

// Lock por cuenta para transferencias concurrentes (se toma en orden de id).
public final ReentrantLock lock = new ReentrantLock();

/** Saldo como decimal exacto (centavos / 100) para el JSON de salida. */
public BigDecimal saldoDecimal() {
    return BigDecimal.valueOf(saldoCentavos, 2);
}
\end{lstlisting}

Cada transferencia confirmada se representa como \codigo{model/Transaccion.java},
un \codigo{record} inmutable con un numero de secuencia que impone un orden total
sobre las operaciones; ese \codigo{seq} es la base de la replicacion y del log
durable.

\begin{lstlisting}[language=Java,caption={Transaccion con orden total (Transaccion.java)}]
/** Una transferencia registrada con su numero de secuencia (orden total). */
public record Transaccion(long seq, int origenId, int destinoId, long montoCentavos) {
}
\end{lstlisting}

El almacen es \codigo{data/CuentaRepository.java}, un singleton sin SGBD que
guarda las cuentas en un \codigo{ConcurrentHashMap} y lleva la secuencia con un
\codigo{AtomicLong}. Su metodo \codigo{transferir()} es atomico: valida monto y
cuentas, bloquea las dos cuentas en orden de \codigo{id}, mueve el saldo, asigna
el \codigo{seq} y notifica al observador del log durable.

\begin{lstlisting}[language=Java,caption={Bloqueo ordenado y conmutacion atomica del saldo (CuentaRepository.java)}]
// Bloqueo en orden de id (menor primero) para no provocar deadlock.
Cuenta primero = origenId < destinoId ? origen : destino;
Cuenta segundo = origenId < destinoId ? destino : origen;
primero.lock.lock();
segundo.lock.lock();
try {
    if (origen.getSaldoCentavos() < montoCentavos) {
        return Resultado.de(Estado.SALDO_INSUFICIENTE);
    }
    origen.setSaldoCentavos(origen.getSaldoCentavos() - montoCentavos);
    destino.setSaldoCentavos(destino.getSaldoCentavos() + montoCentavos);
    long seq = secuencia.incrementAndGet();
    // ...
}
\end{lstlisting}

\begin{advertencia}
No hay base de datos. Todo el estado es volatil: si el proceso del lider
termina, la unica copia durable es el log en Cloud Storage. Las cuentas creadas
en caliente via \codigo{/api/register} que no esten en el CSV no se recuperan tras
un reinicio en frio.
\end{advertencia}

\subsection{Cableado del servidor en BancoHttp}

El servidor HTTP de cada nodo se arma en \codigo{server/BancoHttp.java}, mediante
\codigo{crear()}, compartido por \codigo{MainServer} y por las pruebas. Ahi se
registra cada ruta con su \emph{handler} y se indica en que reactor se ejecuta:
las cuatro operaciones del banco se sirven INLINE en el hilo de I/O salvo
\codigo{/api/login} y \codigo{/api/register} (bcrypt) y \codigo{/panel} (consulta
a los peers), que van al pool WORKER por ser bloqueantes.

\begin{lstlisting}[language=Java,caption={Registro de rutas, handlers y reactores (BancoHttp.java)}]
Ruteador r = new Ruteador()
        .registrar("/api/register", auth, Ruteador.WORKER)
        .registrar("/api/login", auth, Ruteador.WORKER)
        .registrar("/api/accounts", new AccountHandler(), Ruteador.INLINE)
        .registrar("/api/transactions/transfer", new TransferHandler(), Ruteador.INLINE)
        .registrar("/stats", new StatsHandler(), Ruteador.INLINE)
        .registrar("/panel", new PanelHandler(idLocal, peers), Ruteador.WORKER)
        .registrar("/", new PaginaHandler(), Ruteador.INLINE);
return new ServidorNio(puerto, r, reactores(), hilosWorker());
\end{lstlisting}

El numero de reactores (hilos de I/O) y de hilos worker se calcula a partir de
los nucleos disponibles y se puede ajustar con \codigo{BANCO\_REACTORES} y
\codigo{BANCO\_WORKERS}; en una maquina de 4 vCPU se usan 3 reactores por defecto.

\subsubsection{Flujo de alto nivel de una peticion}

Una peticion al lider recorre las siguientes etapas:

\begin{itemize}
\item \codigo{MainServer} carga el CSV con \codigo{AlumnosLoader.cargar()} y crea
  el servidor con \codigo{BancoHttp.crear()}, que devuelve un \codigo{ServidorNio}.
\item El \codigo{ServidorNio} acepta la conexion en un reactor; el
  \codigo{Ruteador} elige el handler por el prefijo de la ruta y lo ejecuta INLINE
  o en el pool WORKER segun lo registrado.
\item El handler (p. ej. \codigo{TransferHandler}) parsea el cuerpo con
  \codigo{json/Json.java}, opera sobre \codigo{CuentaRepository.get()} y, en el
  caso de transferir, dispara \codigo{transferir()}, que asigna el \codigo{seq} y
  notifica al observador del log durable.
\item La respuesta vuelve serializada como JSON con \codigo{Json.toJson()}.
\end{itemize}

El parseo y la serializacion JSON se centralizan en \codigo{json/Json.java}, que
envuelve a gson:

\begin{lstlisting}[language=Java,caption={Fachada JSON centralizada (Json.java)}]
public static JsonObject parse(String body) {
    return JsonParser.parseString(body).getAsJsonObject();
}

public static String toJson(Object o) {
    return GSON.toJson(o);
}
\end{lstlisting}

\subsection{Estructura de paquetes}

El codigo vive bajo \codigo{com.escom.banco} con esta separacion de
responsabilidades:

\begin{tablaglab}{L{4.5cm} L{8.5cm}}
\thead{Paquete} & \thead{Responsabilidad} \\ \midrule
model       & Entidades de dominio: \codigo{Cuenta} y \codigo{Transaccion}. \\
data        & Almacen en memoria y carga: \codigo{CuentaRepository}, \codigo{AlumnosLoader}, \codigo{RegistroTransacciones}. \\
server      & Nodo HTTP y handlers de las rutas (\codigo{MainServer}, \codigo{BancoHttp}, los \codigo{*Handler}). \\
server.http & Capa HTTP propia: \codigo{Ruteador}, \codigo{ParserHttp}, \codigo{Solicitud}, \codigo{Respuesta}. \\
server.nio  & Servidor no bloqueante: \codigo{ServidorNio}, \codigo{Reactor}, \codigo{Conexion}. \\
replicacion & Lider/replica por TCP: \codigo{ServidorReplicacion}, \codigo{ClienteReplicacion}, \codigo{CodecTransaccion}. \\
almacen     & Log durable en Cloud Storage: \codigo{AlmacenGcs}, \codigo{GcsSubidor}. \\
auth        & Identidad y JWT: \codigo{AuthService}, \codigo{JWTUtil}, \codigo{PasswordUtil}. \\
json        & Fachada de serializacion sobre gson: \codigo{Json}. \\
carga       & Generador de carga para pruebas: \codigo{GeneradorCarga}. \\
\end{tablaglab}
